\chapter{Localization of the response and of correlation derivatives}\label{resp:chapter}

In this section, we study complete rows of the response and of both partial derivatives of the correlation and show that they converge to their stationary limits. Our main result, \cref{resp:strong-row}, states that for every $p\ge3$ and $\beta>0$, the lag row $v\mapsto R(t,t-v)$, extended by zero, converges to $r$ in $L^2(0,\infty)$. From it we deduce uniform convergence of response rows (\cref{resp:uniform-row}). We also show that the earlier derivative $\partial_sC(t,s)$ converges in the same $L^2$ sense, so that the squared backward FDT defect vanishes on whole rows (\cref{resp:backward-fdt}). Further, we bound the later derivative $\partial_tC(t,s)$ far from the diagonal for every $p\ge3$ (\cref{resp:forward-derivative}), prove a smoothing estimate in the earlier time (\cref{resp:remote-earlier}), and prove that the thermal part of the Green covariance identity converges uniformly on late triangles (\cref{resp:thermal}).

We use structural identities from \cref{sec:ck-structure} (\cref{prop:green-covariance,prop:uniform-bounds,prop:backward-response,prop:response-compactness}), the Fisher bound of \cref{thm:fisher-clock}, convergence and forgetting from \cref{sec:stationary-selection} (\cref{sel:all-temperature,sel:fixed-wait-forgetting}), and the causal tree realization (\cref{prop:geometry-tree-realization,lem:geometry-projected-drift,sel:prefix-reflection,sel:cluster-complete-monotonicity}). When the plateau is positive, the proof of \cref{resp:strong-row} also uses the cubic tree ceiling of \cref{sec:tree-ceilings} if $p=3$, and the finite forks of \cref{app:general-selection} for general $p$. We give the second argument at the end of this section, in \cref{resp:appendix-p}. The results of this section are used in \cref{resp:integrated-chapter}, which relies on \cref{resp:backward-fdt}, and in \cref{sec:three-time,resp:fisher-chapter}, which use \cref{resp:strong-row,resp:uniform-row,resp:forward-derivative}, together with the Fisher bound \eqref{resp:fisher-domination}. Combined with \cref{thm:fisher-clock}, they prove \eqref{eq:main-strong-rows}, \eqref{eq:main-uniform-response}, and \eqref{eq:main-defect-l2} in \cref{thm:main-response}.

All convergence statements in this section are in the $L^2$ norm or in the uniform norm on a row. They do not control the integral of a row over its remote part. In \cref{resp:integrated-chapter}, we prove that the backward FDT defect, although small in $L^2$, keeps a positive integrated mass there.

\paragraph{Standing assumptions and notation.}
Throughout this section, an integer $p\ge3$ and $\beta>0$ are fixed, except where a statement restricts $p$. Let $(c,r)$ be the stationary profile of \cref{sel:all-temperature}, and let $q=\lim_{v\to\infty}c(v)\in[0,1)$ be its plateau. Recall from \cref{prop:stationary-clusters,sel:cluster-complete-monotonicity} that
\[
 c(0)=r(0)=1,\qquad r=-c',\qquad \int_0^\infty r(v)\,\dd v=1-q,
\]
and that $c$ and $r$ are completely monotone on $(0,\infty)$. In particular, $r$ is nonnegative and nonincreasing, and $r(v)\to0$ as $v\to\infty$ because $r$ is integrable.

By \cref{sel:all-temperature}, the translated kernels converge along every sequence of times tending to infinity. Taking the translation time equal to $t-A$, we find that, for every fixed $0\le A<\infty$,
\begin{equation}\label{eq:resp-recent-strip}
 \lim_{t\to\infty}\sup_{0\le v\le A}\bigl\{|C(t,t-v)-c(v)|+|R(t,t-v)-r(v)|\bigr\}=0.
\end{equation}
We call $\{(t,s):t-A\le s\le t\}$ the recent strip of length $A$. Recall also the forgetting statement \eqref{sel:fixed-wait-decay}: for every fixed $S\ge0$, we have $\sup_{0\le s\le S}\{C(t,s)+R(t,s)\}\to0$ as $t\to\infty$.

We use the bounds $0\le R(t,s)\le C(t,s)\le1$ for $t\ge s$ from \cref{prop:response-compactness}, and $1\le\lambda\le\Lambda$ from \eqref{eq:uniform-multiplier}. Since the memory integral in the response equation \eqref{eq:ck-response} is nonnegative, we have $\partial_tR(t,s)\ge-\lambda(t)R(t,s)$ for $t>s$. Integrating this inequality from the diagonal, where $R(s,s)=1$, we find
\begin{equation}\label{eq:resp-response-lower}
 R(t,s)\ge\exp\Bigl(-\int_s^t\lambda(u)\,\dd u\Bigr)\ge e^{-\Lambda(t-s)},\qquad t\ge s\ge0.
\end{equation}

For $t\ge0$, we extend the lag row of the response by zero:
\begin{equation}
 f_t(v)=\mathbf1_{[0,t]}(v)R(t,t-v),\qquad v\ge0.
 \label{resp:lag-row}
\end{equation}
We write $\norm{\cdot}_2$ for the norm of $L^2(0,\infty)$. By \eqref{eq:initial-response-row-bound}, we have $\norm{f_t}_2^2=\int_0^tR(t,u)^2\,\dd u\le1/2$. For $t\ge L\ge0$, we write
\[
 S_2(t,L)=\int_0^{t-L}R(t,u)^2\,\dd u=\int_L^\infty f_t(v)^2\,\dd v
\]
for the squared response in the remote part of the row. Here and below, an integral with upper endpoint $t-L$ is considered only for $t\ge L$.

Let $v_t$ be the tree vectors of \cref{prop:geometry-tree-realization}, so that $\ip{v_t}{v_s}=C(t,s)$ and $\norm{v_t}=1$. Let $P_T$ be the prefix projections, let $\Pi_j$ be the projections onto trees of height $j$, and let $h(v)=\sqrt p\ip{v}{\mathcal I(v^{\otimes(p-1)})}$ be the tree energy. Recall that $v_s$ belongs to the range of $P_T$ when $s\le T$. For $T,a\ge0$, we set $M_T(a)=P_Tv_{T+a}$. As the whole translated family converges, \cref{sel:prefix-reflection} applies along every sequence $T\to\infty$. For every fixed $L\ge0$, it implies that
\begin{equation}
 \norm{M_T(L)}^2\longrightarrow d_L=c(2L),\qquad
 h(M_T(L))\longrightarrow h_L\qquad(T\to\infty),
 \label{resp:prefix-norm}
\end{equation}
where $h_L$ is given by \eqref{sel:prefix-energy} with $e=e_\beta(\infty)$. If $q>0$, then $d_L\ge q>0$, and we define
\[
 \delta_L=\Eth-\frac{h_L}{d_L^{p/2}}.
\]
By \cref{prop:threshold-inputs}\ref{item:threshold-ceiling}, we have $\delta_L\ge0$ for every $L$. Recall also that the plateau is then marginal, by \cref{sel:positive-marginality}. We prove in \eqref{resp:p-center-deficit} below that $\delta_L\to0$ as $L\to\infty$ for every $p\ge3$. For $p\in\{3,4\}$, both properties are part of \cref{sel:threshold-centers}, whose hypotheses hold by \cref{sel:all-temperature,sel:positive-marginality}. In words, the tree energy of the normalized prefix centers $M_T(L)/\norm{M_T(L)}$ approaches the threshold energy $\Eth$ when $T\to\infty$ first and then $L\to\infty$.

\section{The full response row}\label{sec:resp-full-row}

\begin{theorem}[Strong localization of the response]
\label{resp:strong-row}
Fix $p\ge3$ and $\beta>0$. Then
\begin{equation}
 \lim_{t\to\infty}\norm{f_t-r}_{L^2(0,\infty)}=0.
 \label{resp:strong-row-limit}
\end{equation}
Further, for each fixed $L\ge0$,
\begin{equation}
 \lim_{t\to\infty}\int_0^{t-L}R(t,u)^2\,\dd u=\int_L^\infty r(v)^2\,\dd v,
 \qquad \lim_{t\to\infty}\frac{\chi(t,0)}{\sqrt t}=0.
 \label{resp:square-tail}
\end{equation}
\end{theorem}

The main step is to show that the response keeps no squared mass in the remote part of the row:
\begin{equation}\label{eq:resp-remote-vanishing}
 \lim_{L\to\infty}\limsup_{t\to\infty}S_2(t,L)=0.
\end{equation}
Once \eqref{eq:resp-remote-vanishing} is known, the theorem follows from stationary convergence on the recent strip. We prove \eqref{eq:resp-remote-vanishing} in three cases, which use different inputs.
\begin{itemize}
\item \emph{$q=0$ and every $p\ge3$.} The white-leaf component of the prefix center $M_T(L)$ carries the squared response $2S_2(T+L,L)$, and its norm is at most $\norm{M_T(L)}$. Here we use only the tree realization (\cref{prop:geometry-tree-realization}) and reflection (\cref{sel:prefix-reflection}).
\item \emph{$q>0$ and $p=3$.} We normalize the prefix center and control its white-leaf component by the height inequality of the cubic tree ceiling (\cref{thm:geometry-tree-ceiling}). Besides the inputs of the first case, we need that $\delta_L\to0$, which is part of \cref{sel:threshold-centers} in \cref{sec:tree-ceilings}.
\item \emph{$q>0$ and every $p\ge3$.} We test the mean of finitely many independent continuations against a Gaussian direction built from the Brownian path, and use the overlap penalty of \cref{sel:absolute-field-cost} for an independent direction. This case relies on the finite forks of \cref{app:general-selection} and is carried out in \cref{resp:appendix-p} (\cref{prop:remote-response-general}).
\end{itemize}
The third case contains the second. We include the second because it is short and uses only the deterministic tree geometry.

\begin{proof}[Proof of \cref{resp:strong-row}]
We prove \eqref{eq:resp-remote-vanishing} in Steps~1--3 and deduce the theorem from it in Step~4.

\paragraph{Step 1: the white-leaf bound.}
Fix $T,L\ge0$. By \eqref{eq:geometry-tree-realization}, the component of $v_{T+L}$ in the one-leaf summand $\mathbb Re_0\oplus L^2(\R_+)$ is $R(T+L,0)e_0+\sqrt2\,\mathbf1_{[0,T+L]}R(T+L,\cdot)$. The remaining term of \eqref{eq:geometry-tree-realization} lies in the range of the grafting map, which is orthogonal to that summand. Recall that $P_T$ commutes with $\Pi_0$, and that on the one-leaf summand it restricts the noise-time function to $[0,T]$. We conclude that
\begin{equation}
 \norm{\Pi_0M_T(L)}^2=R(T+L,0)^2+2\int_0^TR(T+L,u)^2\,\dd u.
 \label{resp:white-prefix-norm}
\end{equation}
As $\Pi_0$ is an orthogonal projection, \eqref{resp:white-prefix-norm} implies that $2S_2(T+L,L)\le\norm{\Pi_0M_T(L)}^2\le\norm{M_T(L)}^2$. Letting $T\to\infty$ and using \eqref{resp:prefix-norm}, we find that, with $t=T+L$,
\begin{equation}
 \limsup_{t\to\infty}S_2(t,L)\le\frac{d_L}2=\frac{c(2L)}2.
 \label{resp:basic-tail-bound}
\end{equation}
This bound holds for every $p\ge3$ and every plateau. If $q=0$, then $c(2L)\to0$ as $L\to\infty$, and \eqref{resp:basic-tail-bound} proves \eqref{eq:resp-remote-vanishing}.

\paragraph{Step 2: a positive plateau in the cubic model.}
Suppose that $q>0$ and $p=3$, and fix $L\ge0$. For all large $T$, we have $\norm{M_T(L)}^2\ge q/2$ by \eqref{resp:prefix-norm}; for such $T$, let $z_T=M_T(L)/\norm{M_T(L)}$. For $p=3$, the tree energy is $h=\sqrt3f$, where $f(v)=\ip{v}{\mathcal I(v\otimes v)}$ is the functional of \cref{thm:geometry-tree-ceiling}, and $\Eth=2\sqrt{2/3}=\sqrt3F_*$ with $F_*=2\sqrt2/3$. Keeping only the term $j=0$ in the first inequality of \eqref{eq:geometry-height-defect}, we have
\[
 \norm{\Pi_0z_T}^4\le4\sqrt2\,\bigl[F_*-f(z_T)\bigr]=4\sqrt{2/3}\,\bigl[\Eth-h(z_T)\bigr].
\]
On the other hand, \eqref{resp:white-prefix-norm} implies that $\norm{\Pi_0z_T}^2\ge2S_2(T+L,L)/\norm{M_T(L)}^2$. Combining the two bounds, we find
\[
 S_2(T+L,L)\le\frac{\norm{M_T(L)}^2}2\Bigl(4\sqrt{2/3}\,\bigl[\Eth-h(z_T)\bigr]\Bigr)^{1/2}.
\]
The tree energy is homogeneous of degree three, which gives $h(z_T)=h(M_T(L))/\norm{M_T(L)}^3$, and \eqref{resp:prefix-norm} implies that $h(z_T)\to h_L/d_L^{3/2}=\Eth-\delta_L$. Letting $T\to\infty$ with $L$ fixed, we obtain
\begin{equation}
 \limsup_{t\to\infty}S_2(t,L)
 \le\frac{d_L}2\sqrt{4\sqrt{2/3}\,\delta_L}.
 \label{resp:threshold-tail-bound}
\end{equation}
\Cref{sel:threshold-centers} states that $\delta_L\to0$ as $L\to\infty$. As $d_L\le1$, the right side of \eqref{resp:threshold-tail-bound} tends to zero, and \eqref{eq:resp-remote-vanishing} follows.

\paragraph{Step 3: a positive plateau for general $p$.}
Suppose that $q>0$ and $p\ge3$. \Cref{prop:remote-response-general} states that, for every $L>0$,
\[
 \limsup_{t\to\infty}S_2(t,L)\le\frac{d_L}2\,\bd^{-2/p}\,\delta_L^{2/p},
\]
and the right side tends to zero as $L\to\infty$ by \eqref{resp:p-center-deficit}. The proof of \cref{prop:remote-response-general}, in \cref{resp:appendix-p}, is independent of \cref{resp:strong-row} and its consequences. This completes the proof of \eqref{eq:resp-remote-vanishing} in all cases.

\paragraph{Step 4: from the remote bound to strong convergence.}
We first check that $r\in L^2(0,\infty)$. For fixed $V$, we have $f_t\to r$ uniformly on $[0,V]$ by \eqref{eq:resp-recent-strip}, so $\int_0^Vr^2=\lim_t\int_0^Vf_t^2\le1/2$. Letting $V\to\infty$, we find $\norm{r}_2^2\le1/2$.

Next, fix $L>0$ and $t\ge L$. Splitting $\norm{f_t-r}_2^2$ at $L$, and using $(x-y)^2\le2x^2+2y^2$ on $[L,\infty)$ together with $\int_L^\infty f_t^2=S_2(t,L)$, we have
\[
 \norm{f_t-r}_2^2
 \le L\sup_{0\le v\le L}|R(t,t-v)-r(v)|^2+2S_2(t,L)+2\int_L^\infty r(v)^2\,\dd v.
\]
The first term tends to zero by \eqref{eq:resp-recent-strip}, and taking the upper limit in $t$, we find
\[
 \limsup_{t\to\infty}\norm{f_t-r}_2^2\le2\limsup_{t\to\infty}S_2(t,L)+2\int_L^\infty r(v)^2\,\dd v.
\]
Letting $L\to\infty$, and using \eqref{eq:resp-remote-vanishing} and $r\in L^2$, we obtain \eqref{resp:strong-row-limit}.

For the first limit in \eqref{resp:square-tail}, fix $L\ge0$. Using the triangle inequality, and the fact that restriction to $[L,\infty)$ is a contraction of $L^2(0,\infty)$, we have
\[
 \Bigl|S_2(t,L)^{1/2}-\Bigl(\int_L^\infty r(v)^2\,\dd v\Bigr)^{1/2}\Bigr|
 \le\norm{\mathbf1_{[L,\infty)}(f_t-r)}_2\le\norm{f_t-r}_2,
\]
which tends to zero by \eqref{resp:strong-row-limit}. Finally, note that $\chi(t,0)=\int_0^tR(t,u)\,\dd u=\int_0^tf_t(v)\,\dd v$. Using the Cauchy--Schwarz inequality on $[0,t]$, we obtain
\[
 0\le\chi(t,0)\le\int_0^tr(v)\,\dd v+\int_0^t|f_t(v)-r(v)|\,\dd v\le1-q+\sqrt t\,\norm{f_t-r}_2.
\]
After division by $\sqrt t$, the second limit in \eqref{resp:square-tail} follows from \eqref{resp:strong-row-limit}.
\end{proof}

We next upgrade \cref{resp:strong-row} to uniform convergence. The backward response equation prevents $s\mapsto R(t,s)$ from increasing quickly. A large remote value $R(t,s)$ would then force squared mass on the interval $[s-1,s]$, which \eqref{resp:square-tail} excludes.

\begin{corollary}[Uniform absolute localization]
\label{resp:uniform-row}
Fix $p\ge3$ and $\beta>0$. Then
\begin{equation}
 \lim_{t\to\infty}\sup_{0\le s\le t}|R(t,s)-r(t-s)|=0.
 \label{resp:uniform-row-limit}
\end{equation}
In particular, if $s_n\to\infty$ and $t_n-s_n\to\infty$, then $R(t_n,s_n)\to0$.
\end{corollary}

\begin{proof}
Recall that the kernel $\Sigma$ in \eqref{eq:self-energy-kernel} is nonnegative. Since $R\ge0$, the backward response equation \eqref{eq:backward-response} implies that $\partial_sR(t,s)\le\lambda(s)R(t,s)\le\Lambda R(t,s)$ for $t>s$. Then $w\mapsto e^{-\Lambda w}R(t,w)$ is nonincreasing on $[0,t]$, that is, $R(t,u)\ge e^{-\Lambda(s-u)}R(t,s)$ for $u\le s\le t$. Squaring this inequality and integrating it over $u\in[s-1,s]$, we have, for $1\le s\le t$,
\[
 R(t,s)^2\,\frac{1-e^{-2\Lambda}}{2\Lambda}\le\int_{s-1}^sR(t,u)^2\,\dd u.
\]
If moreover $s\le t-L$, the integral on the right is at most $S_2(t,L)$. Using \eqref{resp:square-tail}, we find that, for every fixed $L\ge0$,
\[
 \limsup_{t\to\infty}\sup_{1\le s\le t-L}R(t,s)
 \le\Bigl(\frac{2\Lambda}{1-e^{-2\Lambda}}\int_L^\infty r(v)^2\,\dd v\Bigr)^{1/2}.
\]
Recall from \eqref{sel:fixed-wait-decay} that $\sup_{0\le s\le1}R(t,s)\to0$. On the recent strip $t-L\le s\le t$, the difference $|R(t,s)-r(t-s)|$ tends to zero uniformly by \eqref{eq:resp-recent-strip}. On the remaining part $0\le s\le t-L$, we have $|R(t,s)-r(t-s)|\le R(t,s)+r(L)$, since $r$ is nonincreasing. Combining these three estimates, we have
\[
 \limsup_{t\to\infty}\sup_{0\le s\le t}|R(t,s)-r(t-s)|
 \le\Bigl(\frac{2\Lambda}{1-e^{-2\Lambda}}\int_L^\infty r(v)^2\,\dd v\Bigr)^{1/2}+r(L).
\]
The right side tends to zero as $L\to\infty$, which proves \eqref{resp:uniform-row-limit}. For the last assertion, note that $R(t_n,s_n)\le|R(t_n,s_n)-r(t_n-s_n)|+r(t_n-s_n)$, and both terms tend to zero.
\end{proof}

\section{The later correlation derivative}\label{sec:resp-later-derivative}

We now bound $\partial_tC(t,s)$ when $t-s$ is large. By \cref{lem:geometry-projected-drift}, the projected curve $a\mapsto M_T(a)$ is continuously differentiable for $a>0$, with
\[
 \frac{\dd}{\dd a}M_T(a)=P_T\mathcal B(v_{T+a}),\qquad
 \mathcal B(v)=-v+\beta\nabla h(v)-p\beta h(v)v,
\]
and $\mathcal B$ is Lipschitz on the closed unit ball of $\mathcal H$. Reflection identifies limits of inner products of prefix centers, not of their derivatives. For this reason, we pass through difference quotients. This requires a modulus of continuity for the projected derivative that is uniform in $T$, which we now record.

\begin{lemma}[Hölder continuity of the projected drift]\label{lem:resp-drift-holder}
Fix $p\ge3$ and $\beta>0$. Let $L_{\mathcal B}$ be a Lipschitz constant of $\mathcal B$ on the closed unit ball of $\mathcal H$, and set $\kappa_{\mathcal B}=L_{\mathcal B}\sqrt{2\Lambda}$. Then, for all $t\ge s\ge0$,
\[
 0\le1-C(t,s)\le\Lambda(t-s),\qquad \norm{v_t-v_s}\le\sqrt{2\Lambda(t-s)},
\]
and, for all $t,t'\ge0$ and $T\ge0$,
\[
 \norm{\mathcal B(v_t)-\mathcal B(v_{t'})}\le\kappa_{\mathcal B}|t-t'|^{1/2},
 \qquad
 \norm{P_T\mathcal B(v_t)-P_T\mathcal B(v_{t'})}\le\kappa_{\mathcal B}|t-t'|^{1/2}.
\]
\end{lemma}

\begin{proof}
The first bound is \eqref{eq:diagonal-increment} in \cref{prop:global-regularity}. The second bound follows from the identity $\norm{v_t-v_s}^2=2-2C(t,s)$. The vectors $v_t$ lie in the unit ball, and the third bound follows from the second and the Lipschitz property of $\mathcal B$. The last bound holds because $P_T$ is a contraction.
\end{proof}

\begin{theorem}[Remote derivative bound and forward FDT]
\label{resp:forward-derivative}
Fix an integer $p\ge3$ and $\beta>0$. For every fixed $L>0$,
\begin{equation}
 \limsup_{t\to\infty}\sup_{0\le s\le t-L}
 |\partial_tC(t,s)|\le\sqrt{c''(2L)}
 \le\frac{\sqrt{1-q}}{eL}.
 \label{resp:remote-forward}
\end{equation}
Further,
\begin{equation}
 \lim_{t\to\infty}\sup_{0\le s\le t}|\partial_tC(t,s)-c'(t-s)|=0,
 \qquad
 \lim_{t\to\infty}\sup_{0\le s\le t}|\partial_tC(t,s)+R(t,s)|=0.
 \label{resp:uniform-forward-fdt}
\end{equation}
Here $\partial_tC(t,t)$ is the right derivative of $u\mapsto C(u,t)$ at $u=t$; it equals $-1$ for every $t\ge0$. The value $c'(0)=-1$ is also the right derivative.
\end{theorem}

\begin{proof}
All inputs to this proof hold for every integer $p\ge3$: the projected drift of \cref{lem:geometry-projected-drift}, its uniform modulus in \cref{lem:resp-drift-holder}, prefix reflection in \cref{sel:prefix-reflection}, complete monotonicity in \cref{sel:cluster-complete-monotonicity}, and uniform response localization in \cref{resp:uniform-row}.

The proof has three steps. First, we bound the norm of the projected derivative $M_T'(L)$ through difference quotients and reflection. Since $\partial_tC(T+L,s)=\ip{M_T'(L)}{v_s}$ for every $s\le T$, we deduce \eqref{resp:remote-forward} from this bound. Second, we show that $\partial_tC(t,t-v)\to c'(v)$ uniformly on each recent strip, by differentiating stationary convergence through forward difference quotients. Third, we combine the two estimates with \cref{resp:uniform-row}.

\paragraph{Step 1: the remote bound.}
Fix $L>0$ and $\eta\in(0,1]$. For every $T\ge0$, we have $M_T(L+\eta)-M_T(L)=\int_0^\eta M_T'(L+w)\,\dd w$. \Cref{lem:resp-drift-holder} then implies that
\[
 \Bigl\|M_T'(L)-\frac{M_T(L+\eta)-M_T(L)}\eta\Bigr\|
 \le\frac1\eta\int_0^\eta\norm{M_T'(L+w)-M_T'(L)}\,\dd w
 \le\frac23\kappa_{\mathcal B}\sqrt\eta.
\]
Applying \eqref{sel:prefix-gram} to the three inner products in the expansion of the squared norm, we have
\[
 \lim_{T\to\infty}\Bigl\|\frac{M_T(L+\eta)-M_T(L)}\eta\Bigr\|^2
 =\frac{c(2L+2\eta)-2c(2L+\eta)+c(2L)}{\eta^2}.
\]
Combining the last two displays, we find
\[
 \limsup_{T\to\infty}\norm{M_T'(L)}
 \le\frac23\kappa_{\mathcal B}\sqrt\eta
 +\Bigl(\frac{c(2L+2\eta)-2c(2L+\eta)+c(2L)}{\eta^2}\Bigr)^{1/2}.
\]
The left side does not depend on $\eta$. The function $c$ is smooth on $(0,\infty)$, and the second difference quotient tends to $c''(2L)$ as $\eta\downarrow0$. Taking $T\to\infty$ first and then $\eta\downarrow0$, we have
\[
 \limsup_{T\to\infty}\norm{M_T'(L)}\le\sqrt{c''(2L)}.
\]
If $s\le T$, then $v_s$ lies in the range of $P_T$, so $C(T+a,s)=\ip{M_T(a)}{v_s}$ for all $a\ge0$, and $\partial_tC(T+L,s)=\ip{M_T'(L)}{v_s}$. As $\norm{v_s}=1$, we have $|\partial_tC(T+L,s)|\le\norm{M_T'(L)}$ for every $s\le T$. Setting $T=t-L$ proves the first inequality in \eqref{resp:remote-forward}.

For the second inequality, we write the representation \eqref{eq:stationary-spectral-measure} as $c(v)=q+\int_{(0,\infty)}e^{-xv}\,\nu(\dd x)$, where $\nu((0,\infty))=1-q$. Differentiating under the integral for $v>0$, we have
\[
 c''(2L)=\int_{(0,\infty)}x^2e^{-2Lx}\,\nu(\dd x)\le\frac{1-q}{e^2L^2},
\]
where we used that $\sup_{x>0}x^2e^{-2Lx}=(eL)^{-2}$, attained at $x=1/L$.

\paragraph{Step 2: the recent strip.}
For $t>s$, we apply the projected drift with $T=s$ to $C(t,s)=\ip{P_sv_t}{v_s}$. This gives
\[
 \partial_tC(t,s)=\ip{P_s\mathcal B(v_t)}{v_s}=\ip{\mathcal B(v_t)}{v_s}.
\]
The right side is continuous in $t\ge s$. At $t=s$ it equals $-1$, because $\ip{\nabla h(v)}{v}=ph(v)$ by homogeneity and $\norm{v_s}=1$; this is the derivative from the causal side at the diagonal. \Cref{lem:resp-drift-holder} also implies that $|\partial_tC(t,s)-\partial_tC(t',s)|\le\kappa_{\mathcal B}|t-t'|^{1/2}$ for all $t,t'\ge s$.

Fix $L>0$ and $\eta\in(0,1]$. For $0\le v\le L$ and $t\ge L$, we hold the earlier time $t-v$ fixed in a forward difference quotient. The Hölder bound just stated gives
\[
 \Bigl|\partial_tC(t,t-v)-\frac{C(t+\eta,t-v)-C(t,t-v)}\eta\Bigr|
 \le\frac1\eta\int_0^\eta\kappa_{\mathcal B}\sqrt w\,\dd w=\frac23\kappa_{\mathcal B}\sqrt\eta.
\]
By \eqref{eq:resp-recent-strip} with $A=L+\eta$, the difference quotient converges to $[c(v+\eta)-c(v)]/\eta$ uniformly in $v\in[0,L]$. We obtain
\[
 \limsup_{t\to\infty}\sup_{0\le v\le L}|\partial_tC(t,t-v)-c'(v)|
 \le\frac23\kappa_{\mathcal B}\sqrt\eta
 +\sup_{0\le v\le L}\Bigl|\frac{c(v+\eta)-c(v)}\eta-c'(v)\Bigr|.
\]
The function $c'=-r$ is uniformly continuous on $[0,L+1]$, including the right derivative $c'(0)=-1$, and the last supremum tends to zero as $\eta\downarrow0$. We conclude that $\sup_{t-L\le s\le t}|\partial_tC(t,s)-c'(t-s)|\to0$ for every fixed $L$.

\paragraph{Step 3: the whole row.}
On the remote part $0\le s\le t-L$, we have $|\partial_tC(t,s)-c'(t-s)|\le|\partial_tC(t,s)|+r(L)$. Using Step 2 and \eqref{resp:remote-forward}, we find
\[
 \limsup_{t\to\infty}\sup_{0\le s\le t}|\partial_tC(t,s)-c'(t-s)|\le\sqrt{c''(2L)}+r(L).
\]
Letting $L\to\infty$ proves the first limit in \eqref{resp:uniform-forward-fdt}. As $c'=-r$, the second limit follows from the first and \cref{resp:uniform-row}.
\end{proof}

\section{The earlier derivative and the squared FDT defect}\label{sec:resp-earlier-derivative}

For $t>s>0$, let $E(t,s)=\partial_sC(t,s)-R(t,s)$ be the backward FDT defect. By \cref{thm:fisher-clock}, the derivative $\partial_sC$ is continuous on the open causal triangle, and there exists a nonnegative integrable function $a_F$, obtained on one dimension subsequence, such that
\begin{equation}
 |E(t,s)|^2\le a_F(s)\quad(t>s>0),\qquad
 \int_0^\infty a_F(s)\,\dd s<\infty,\qquad
 \lim_{S\to\infty}\sup_{s\ge S}a_F(s)=0.
 \label{resp:fisher-domination}
\end{equation}
We do not need $a_F$ to be unique. The bound \eqref{resp:fisher-domination} is the spherical form used here of the entropy-production bounds on FDT violation of Cugliandolo, Dean, and Kurchan \cite[arXiv version~2, equations~(12)--(13), p.~3, and~(18), p.~4]{cdk1997}. The common majorant in \eqref{resp:fisher-domination} is proved in \cref{thm:fisher-clock}, and the whole-row conclusions below are proved here. It shows that the defect is uniformly small when the earlier time is late. The row $[0,t]$ grows with $t$, however, and this alone does not control the defect on the whole row. The next theorem shows that the integral of the squared defect over the row nevertheless tends to zero.

\begin{theorem}[Strong backward FDT]
\label{resp:backward-fdt}
Fix $p\ge3$ and $\beta>0$. Then
\begin{equation}
 \lim_{t\to\infty}\sup_{0\le s\le S}|\partial_sC(t,s)|=0
 \quad\text{for each fixed }0\le S<\infty,
 \qquad
 \lim_{t\to\infty}\int_0^t|E(t,s)|^2\,\dd s=0.
 \label{resp:backward-fdt-limit}
\end{equation}
Further, if $g_t(v)=\mathbf1_{[0,t]}(v)\partial_sC(t,t-v)$, then
\begin{equation}
 \lim_{t\to\infty}\norm{g_t-r}_{L^2(0,\infty)}=0.
 \label{resp:earlier-row-limit}
\end{equation}
For $t>0$, the earlier derivative at $s=0$ is the right derivative. At $s=t>0$, it is the left derivative of $u\mapsto C(t,u)$ and equals $1$; we set $E(t,t)=0$.
\end{theorem}

\begin{proof}
We proceed in three steps. First, we represent the colored memory term of the Green identity through feature vectors. A positivity argument, together with forgetting, shows that this term is negligible on every bounded interval of earlier times. Second, we differentiate the Green identity in its earlier time. This expresses $\partial_sC(t,s)$ through $C$, $R$, and the colored term, and it proves the first claim in \eqref{resp:backward-fdt-limit}. Third, the Fisher majorant allows dominated convergence over the whole row, and \cref{resp:strong-row} implies \eqref{resp:earlier-row-limit}.

\paragraph{Step 1: the colored term forgets bounded earlier intervals.}
For $s\ge0$, let $\phi_s=v_s^{\otimes(p-1)}\in\mathcal H^{\otimes(p-1)}$. These are unit vectors with $\ip{\phi_u}{\phi_s}=C(u,s)^{p-1}$ and $\norm{\phi_u-\phi_s}^2=2[1-C(u,s)^{p-1}]$. In particular, $s\mapsto\phi_s$ is norm continuous. We define
\[
 F_t=\int_0^tR(t,u)\phi_u\,\dd u,\qquad
 \mathcal K(t,s)=\ip{F_t}{\phi_s}=\int_0^tR(t,u)C(u,s)^{p-1}\,\dd u,
\]
where the first integral is a Bochner integral of a continuous function on $[0,t]$. Note that $\mathcal K(t,s)\ge0$. When $s=t$, the colored memory term in the Green identity \eqref{eq:green-covariance} equals $\beta^2p\norm{F_t}^2$. Since all three terms of \eqref{eq:green-covariance} are nonnegative, we obtain $\beta^2p\norm{F_t}^2\le C(t,t)=1$. By the Cauchy--Schwarz inequality, we have, for all $t,s,v\ge0$,
\begin{equation}
 |\mathcal K(t,s)-\mathcal K(t,v)|\le(\beta^2p)^{-1/2}\sqrt{2[1-C(s,v)^{p-1}]}.
 \label{resp:colored-modulus}
\end{equation}
This modulus does not depend on $t$. For $s\le t$, Fubini's theorem expresses the colored memory term in \eqref{eq:green-covariance} as $\beta^2p\int_0^sR(s,v)\mathcal K(t,v)\,\dd v$. The other two terms are nonnegative, and therefore
\begin{equation}
 C(t,s)\ge\beta^2p\int_0^sR(s,v)\mathcal K(t,v)\,\dd v.
 \label{resp:colored-test}
\end{equation}

We claim that $\mathcal K(t,s)\to0$ as $t\to\infty$, for every fixed $s>0$. Suppose instead that $\mathcal K(t_n,s)\ge\varepsilon$ along a sequence $t_n\to\infty$, for some $\varepsilon>0$. Since $C$ is continuous and $C(s,s)=1$, there exists $h\in(0,s]$ such that the right side of \eqref{resp:colored-modulus} is at most $\varepsilon/2$ for $v\in[s-h,s]$. Then $\mathcal K(t_n,v)\ge\varepsilon/2$ on $[s-h,s]$. Inserting this and the lower bound \eqref{eq:resp-response-lower} into \eqref{resp:colored-test}, we obtain
\[
 C(t_n,s)\ge\frac{\beta^2p\,\varepsilon}2\int_{s-h}^se^{-\Lambda(s-v)}\,\dd v
 =\frac{\beta^2p\,\varepsilon}{2\Lambda}\bigl(1-e^{-\Lambda h}\bigr)>0,
\]
which contradicts $C(t_n,s)\to0$ from \eqref{sel:fixed-wait-decay}. For $s=0$ and any $s'>0$, \eqref{resp:colored-modulus} implies that $\mathcal K(t,0)\le\mathcal K(t,s')+(\beta^2p)^{-1/2}\sqrt{2[1-C(s',0)^{p-1}]}$. Letting $t\to\infty$ and then $s'\downarrow0$, we find $\mathcal K(t,0)\to0$.

Finally, fix $S$ and $\varepsilon>0$. The correlation is uniformly continuous on $[0,S]^2$ with $C(s,s)=1$. Hence there exists a finite set $\{s_1,\dots,s_k\}\subset[0,S]$ such that every $s\in[0,S]$ has some $s_i$ with $(\beta^2p)^{-1/2}\sqrt{2[1-C(s,s_i)^{p-1}]}\le\varepsilon$. Then \eqref{resp:colored-modulus} implies that $\sup_{0\le s\le S}\mathcal K(t,s)\le\max_i\mathcal K(t,s_i)+\varepsilon$. Letting $t\to\infty$ and then $\varepsilon\downarrow0$, we conclude that
\begin{equation}\label{eq:resp-colored-forgetting}
 \lim_{t\to\infty}\sup_{0\le s\le S}\mathcal K(t,s)=0.
\end{equation}

\paragraph{Step 2: the backward correlation equation.}
Fix $t>0$. By the Green identity \eqref{eq:green-covariance} and Fubini's theorem, we have, for $0\le s\le t$,
\[
 C(t,s)=R(t,0)R(s,0)+2\int_0^sR(t,u)R(s,u)\,\dd u+\beta^2p\int_0^sR(s,v)\mathcal K(t,v)\,\dd v.
\]
By \cref{prop:backward-response}, the derivative $\partial_1R$ is continuous on each compact causal triangle, and $\mathcal K(t,\cdot)$ is continuous. We may therefore differentiate in $s\in(0,t)$ under the integral signs. Using $R(s,s)=1$, we obtain
\begin{align*}
 \partial_sC(t,s)={}&2R(t,s)+\beta^2p\,\mathcal K(t,s)+R(t,0)\,\partial_1R(s,0)\\
 &+2\int_0^sR(t,u)\,\partial_1R(s,u)\,\dd u+\beta^2p\int_0^s\partial_1R(s,v)\mathcal K(t,v)\,\dd v.
\end{align*}
The first two terms are the boundary terms of the thermal and colored integrals. We now insert the response equation \eqref{eq:ck-response}, which reads $\partial_1R(s,v)=-\lambda(s)R(s,v)+\int_v^s\Sigma(s,w)R(w,v)\,\dd w$ with $\Sigma(s,w)=\beta^2p(p-1)C(s,w)^{p-2}R(s,w)$. The terms with $-\lambda(s)$ reproduce $-\lambda(s)$ times the Green representation of $C(t,s)$. In the memory terms, we exchange the order of integration on the triangle $\{0\le u\le w\le s\}$, and they become
\[
 \int_0^s\Sigma(s,w)\Bigl[R(t,0)R(w,0)+2\int_0^wR(t,u)R(w,u)\,\dd u+\beta^2p\int_0^wR(w,v)\mathcal K(t,v)\,\dd v\Bigr]\dd w.
\]
By the Green identity at $(t,w)$, the bracket equals $C(t,w)$. We have shown that, for $0<s<t$,
\begin{align}
 \partial_sC(t,s)={}&-\lambda(s)C(t,s)+2R(t,s)
       +\beta^2p(p-1)\int_0^s
          C(s,u)^{p-2}R(s,u)C(t,u)\,\dd u\notag\\
       &+\beta^2p\,\mathcal K(t,s).
 \label{resp:backward-correlation}
\end{align}
All integrations and differentiations here take place on finite rectangles, where the integrands are continuous and bounded. The right side of \eqref{resp:backward-correlation} is continuous on $[0,t]$, so it also gives the right derivative at $s=0$ and the left derivative at $s=t$.

As a check, we let $s\uparrow t$. Using the definition of $\mathcal K$ and \eqref{eq:ck-multiplier}, we find
\[
\begin{aligned}
 \beta^2p\,\mathcal K(t,t)+\beta^2p(p-1)\int_0^tC(t,u)^{p-1}R(t,u)\,\dd u
 &=\beta^2p^2\int_0^tC(t,u)^{p-1}R(t,u)\,\dd u\\
 &=\lambda(t)-1.
\end{aligned}
\]
Inserting this into \eqref{resp:backward-correlation}, we obtain $\partial_sC(t,t-)=-\lambda(t)+2+\lambda(t)-1=1$, and $E(t,t-)=0$, as the unit diagonal response requires.

Now fix $S$ and let $t>S$. Bounding $C(s,u)^{p-2}\le1$ and $R(s,u)\le1$ in \eqref{resp:backward-correlation}, we have, for $0\le s\le S$,
\[
 |\partial_sC(t,s)|\le\bigl(\Lambda+\beta^2p(p-1)S\bigr)\sup_{0\le u\le S}C(t,u)
 +2\sup_{0\le u\le S}R(t,u)+\beta^2p\sup_{0\le u\le S}\mathcal K(t,u).
\]
Each supremum tends to zero, by \eqref{sel:fixed-wait-decay} and \eqref{eq:resp-colored-forgetting}. This proves the first claim in \eqref{resp:backward-fdt-limit}.

\paragraph{Step 3: the whole row.}
Let $\widetilde E_t(s)=\mathbf1_{(0,t)}(s)E(t,s)$ for $s>0$. For fixed $s>0$, the first claim in \eqref{resp:backward-fdt-limit} and \eqref{sel:fixed-wait-decay} imply that $\widetilde E_t(s)\to0$ as $t\to\infty$. Further, $|\widetilde E_t(s)|^2\le a_F(s)$ for all $t$ by \eqref{resp:fisher-domination}, and $a_F$ is integrable on $(0,\infty)$. By the dominated convergence theorem, we conclude that
\[
 \lim_{t\to\infty}\int_0^t|E(t,s)|^2\,\dd s=\lim_{t\to\infty}\int_0^\infty|\widetilde E_t(s)|^2\,\dd s=0.
\]
This is the second claim in \eqref{resp:backward-fdt-limit}. After the change of variables $s=t-v$, we have $\norm{g_t-f_t}_2^2=\int_0^t|E(t,s)|^2\,\dd s$. Then \cref{resp:strong-row} implies that $\norm{g_t-r}_2\le\norm{g_t-f_t}_2+\norm{f_t-r}_2\to0$, which is \eqref{resp:earlier-row-limit}.
\end{proof}

The proof does not require a unique limit of the Fisher information, nor a majorant chosen anew for each terminal time: one function $a_F$ dominates every row. The positivity of $C$ and $R$ enters through \eqref{resp:colored-test}.

The next corollary states that remote values of the response and of the earlier derivative are small, uniformly on the late triangle, and that remote correlations change little across a window of fixed length in the earlier time. It is used in \cref{sec:three-time}.

\begin{corollary}[Uniform remote decay and earlier smoothing]
\label{resp:remote-earlier}
Fix $p\ge3$ and $\beta>0$. Then
\begin{equation}
 \lim_{D\to\infty}\limsup_{S\to\infty}
 \sup_{s\ge S,\ t\ge s+D}
 \bigl(R(t,s)+|\partial_sC(t,s)|\bigr)=0.
 \label{resp:remote-earlier-limit}
\end{equation}
More precisely, for each fixed $D>0$,
\begin{equation}
 \limsup_{S\to\infty}
 \sup_{s\ge S+D,\ t\ge s+D}
 \int_{s-D}^s|\partial_uC(t,u)|\,\dd u
 \le\sqrt D\left(\int_D^\infty r(v)^2\,\dd v\right)^{1/2}.
 \label{resp:earlier-smoothing}
\end{equation}
If $r(v)=O(v^{-1-a})$ as $v\to\infty$ for some $a>0$, the right side of \eqref{resp:earlier-smoothing} is $O(D^{-a})$ as $D\to\infty$.
\end{corollary}

When $\beta\ge\bd$, \cref{thm:stationary-critical-amplitude} gives $r(v)\sim aA_\beta v^{-1-a}$, with the exponent $a\in(0,1/2)$ defined there, and the last hypothesis holds.

\begin{proof}
Let
\[
 \epsilon(U)=\sup_{t\ge U}\norm{f_t-r}_2,\qquad
 \epsilon_\infty(U)=\sup_{t\ge U}\sup_{0\le s\le t}|R(t,s)-r(t-s)|.
\]
Both tend to zero as $U\to\infty$, by \cref{resp:strong-row,resp:uniform-row}. For $t>s>0$, we have $|\partial_sC(t,s)|\le R(t,s)+\sqrt{a_F(s)}$ by \eqref{resp:fisher-domination}.

For the pointwise statement, let $s\ge S$ and $t\ge s+D$. As $r$ is nonincreasing, $R(t,s)\le r(D)+\epsilon_\infty(S+D)$, which implies
\[
 R(t,s)+|\partial_sC(t,s)|\le2r(D)+2\epsilon_\infty(S+D)+\sup_{u\ge S}\sqrt{a_F(u)}.
\]
We take $S\to\infty$, using \eqref{resp:fisher-domination}, and then $D\to\infty$. This proves \eqref{resp:remote-earlier-limit}.

For the integrated statement, fix $D>0$, and let $s\ge S+D$ and $t\ge s+D$. Then $[s-D,s]\subset[S,t-D]$ and $t\ge S+2D$. Recall that $\mu([S,\infty))=\int_S^\infty a_F(u)\,\dd u$, where $\mu$ is the measure of \eqref{eq:entropy-measure}. By the Cauchy--Schwarz inequality,
\[
\begin{aligned}
 \int_{s-D}^s|\partial_uC(t,u)|\,\dd u
 &\le\sqrt D\Bigl(\int_{s-D}^sR(t,u)^2\,\dd u\Bigr)^{1/2}+\sqrt{D\,\mu([S,\infty))}\\
 &\le\sqrt D\,S_2(t,D)^{1/2}+\sqrt{D\,\mu([S,\infty))}.
\end{aligned}
\]
As in Step 4 of the proof of \cref{resp:strong-row}, we have $S_2(t,D)^{1/2}\le(\int_D^\infty r^2)^{1/2}+\norm{f_t-r}_2$. The supremum in \eqref{resp:earlier-smoothing} is therefore at most
\[
 \sqrt D\Bigl[\Bigl(\int_D^\infty r(v)^2\,\dd v\Bigr)^{1/2}+\epsilon(S+2D)\Bigr]+\sqrt{D\,\mu([S,\infty))}.
\]
Since $a_F$ is integrable, $\mu([S,\infty))\to0$, and letting $S\to\infty$ proves \eqref{resp:earlier-smoothing}. Finally, if $r(v)\le A_0v^{-1-a}$ for $v\ge v_0$, then for $D\ge v_0$ we have
\[
 \sqrt D\Bigl(\int_D^\infty r(v)^2\,\dd v\Bigr)^{1/2}\le\sqrt D\Bigl(\frac{A_0^2D^{-1-2a}}{1+2a}\Bigr)^{1/2}=\frac{A_0}{\sqrt{1+2a}}D^{-a}.
\]
This completes the proof.
\end{proof}

\section{The direct thermal covariance}\label{sec:resp-thermal}

The Green identity \eqref{eq:green-covariance} splits $C(t,s)$ into an initial term, a thermal term, which comes directly from the white noise, and a colored memory term. We define the thermal term and its stationary counterpart by
\[
 B(t,s)=2\int_0^{t\wedge s}R(t,u)R(s,u)\,\dd u,
 \qquad b_{\rm th}(\tau)=2\int_0^\infty r(\tau+v)r(v)\,\dd v.
\]
We say that a sequence of pairs $(t_n,s_n)$ with $t_n\ge s_n$ is a late pair with diverging lag if $s_n\to\infty$ and $t_n-s_n\to\infty$.

\begin{corollary}[Thermal covariance on the full late triangle]
\label{resp:thermal}
Fix $p\ge3$ and $\beta>0$. Then
\begin{equation}
 \lim_{S\to\infty}\sup_{t\ge s\ge S}|B(t,s)-b_{\rm th}(t-s)|=0.
 \label{resp:thermal-uniform}
\end{equation}
The function $b_{\rm th}$ is completely monotone on $(0,\infty)$ and tends to zero at infinity. In particular, $B(t_n,s_n)\to0$ along every late pair with diverging lag, and $B(t,t)\to2\int_0^\infty r(v)^2\,\dd v$. If $C(t_n,s_n)\to c_0$ along a late pair with diverging lag, then
\begin{equation}
 \lim_{n\to\infty}\beta^2p\int_0^{t_n}\int_0^{s_n}
 R(t_n,u)R(s_n,v)C(u,v)^{p-1}\,\dd v\,\dd u=c_0.
 \label{resp:colored-aging}
\end{equation}
If $r(\tau)\sim aA\tau^{-1-a}$ as $\tau\to\infty$ for some constants $a,A>0$, then
\begin{equation}
 b_{\rm th}(\tau)\sim2(1-q)\,r(\tau)
 \sim2(1-q)aA\tau^{-1-a}.
 \label{resp:thermal-tail}
\end{equation}
\end{corollary}

By \cref{thm:stationary-critical-amplitude}, the hypothesis of the last assertion holds when $\beta\ge\bd$, with $A=A_\beta$. The identity \eqref{resp:colored-aging} says that any correlation retained along a late pair with diverging lag comes from the colored memory term, with its whole integration rectangle.

\begin{proof}
Let $\tau=t-s\ge0$. After the substitution $u=s-v$, the zero extension \eqref{resp:lag-row} gives
\[
 B(t,s)=2\int_0^\infty f_t(\tau+v)f_s(v)\,\dd v,
\]
because $f_s(v)=0$ for $v>s$. Let $(\mathcal T_\tau g)(v)=g(\tau+v)$; this translation is a contraction of $L^2(0,\infty)$. Then $B(t,s)=2\ip{\mathcal T_\tau f_t}{f_s}$ and $b_{\rm th}(\tau)=2\ip{\mathcal T_\tau r}{r}$. Using $\norm{f_s}_2,\norm{r}_2\le1/\sqrt2$, we have
\[
 |B(t,s)-b_{\rm th}(\tau)|
 \le2\bigl|\ip{\mathcal T_\tau(f_t-r)}{f_s}\bigr|+2\bigl|\ip{\mathcal T_\tau r}{f_s-r}\bigr|
 \le\sqrt2\,\bigl(\norm{f_t-r}_2+\norm{f_s-r}_2\bigr),
\]
uniformly in $\tau\ge0$. With $\epsilon(S)=\sup_{t\ge S}\norm{f_t-r}_2$, the left side of \eqref{resp:thermal-uniform} is at most $2\sqrt2\,\epsilon(S)$, which tends to zero by \cref{resp:strong-row}.

Next, we write $c(v)=q+\int_{(0,\infty)}e^{-xv}\nu(\dd x)$ as in \eqref{eq:stationary-spectral-measure}. Let $\sigma(\dd x)=x\,\nu(\dd x)$ on $(0,\infty)$. Then $r(v)=\int_{(0,\infty)}e^{-xv}\sigma(\dd x)$ for $v>0$, and $\sigma((0,\infty))=\lim_{v\downarrow0}r(v)=1$ by monotone convergence. By Tonelli's theorem,
\[
 b_{\rm th}(\tau)
 =2\int_{(0,\infty)^2}\frac{e^{-x\tau}}{x+y}\,\sigma(\dd x)\,\sigma(\dd y)
 =\int_{(0,\infty)}e^{-x\tau}\rho_{\rm th}(\dd x),
\]
where $\rho_{\rm th}(\dd x)=2\sigma(\dd x)\int_{(0,\infty)}(x+y)^{-1}\sigma(\dd y)$. The measure $\rho_{\rm th}$ is positive with total mass $b_{\rm th}(0)=2\norm{r}_2^2\le1$. Hence $b_{\rm th}$ is the Laplace transform of a finite positive measure on $(0,\infty)$, and it is completely monotone on $(0,\infty)$. Also $b_{\rm th}(\tau)\le2\norm{r}_2\,\norm{\mathbf1_{[\tau,\infty)}r}_2\to0$.

Along a late pair with diverging lag, we have $|B(t_n,s_n)|\le|B(t_n,s_n)-b_{\rm th}(t_n-s_n)|+b_{\rm th}(t_n-s_n)\to0$. Taking $\tau=0$ in \eqref{resp:thermal-uniform}, we find $B(t,t)\to b_{\rm th}(0)=2\norm{r}_2^2$. The initial term of the Green identity satisfies $0\le R(t_n,0)R(s_n,0)\le R(t_n,0)\to0$, by \eqref{sel:fixed-wait-decay}. Subtracting the initial and thermal terms from \eqref{eq:green-covariance} at $(t_n,s_n)$ proves \eqref{resp:colored-aging}.

For \eqref{resp:thermal-tail}, note that $r>0$ by \cref{prop:scalar-clusters}, and that
\[
 \frac{b_{\rm th}(\tau)}{r(\tau)}=2\int_0^\infty\frac{r(\tau+v)}{r(\tau)}\,r(v)\,\dd v.
\]
As $r$ is nonincreasing, the ratio $r(\tau+v)/r(\tau)$ lies in $[0,1]$. It tends to one for each fixed $v$, by the assumed asymptotics. The measure $r(v)\,\dd v$ has finite mass $1-q$. Using dominated convergence, we find $b_{\rm th}(\tau)/r(\tau)\to2(1-q)$.
\end{proof}

\begin{remark}[Square integrable versus integrated response]\label{rem:resp-l2-versus-l1}
The convergence statements of this section concern $L^2$ norms and uniform norms of rows. A slowly varying contribution spread over a growing part of the row can have vanishing squared norm and a positive integral. This happens for the response in the cubic model. For $p=3$ and $\beta\ge\sqrt{4/3}$, \cref{thm:geometry-susceptibility-lower} gives $\liminf_t\chi(t,0)\ge2(1-q)$, while $\int_0^Lf_t(v)\,\dd v\to\int_0^Lr(v)\,\dd v\le1-q$ for each fixed $L$. We conclude that, for every fixed $L$,
\[
 \liminf_{t\to\infty}\int_0^{t-L}R(t,u)\,\dd u\ge1-q>0,
\]
although $\int_0^{t-L}R(t,u)^2\,\dd u\to\int_L^\infty r^2$, which tends to zero as $L\to\infty$. In \cref{resp:integrated-chapter}, we prove a statement of this kind for the backward FDT defect.
\end{remark}

\section{The remote response for general pure models}
\label{resp:appendix-p}

In this section, we prove \cref{prop:remote-response-general}. It is the case $q>0$ of the remote estimate \eqref{eq:resp-remote-vanishing} for general $p\ge3$, which we used in Step~3 of the proof of \cref{resp:strong-row}. For $p=3$, we controlled the white-leaf component of a normalized prefix center directly, by the height inequality of the cubic tree ceiling. For general $p$, we have no such inequality. Instead, we test the mean of finitely many independent continuations of the dynamics against one Gaussian direction built from the Brownian path. The overlap penalty for a direction independent of the disorder (\cref{sel:absolute-field-cost}) then plays the role of the height inequality. That lemma rests on the external-field algorithmic threshold of Huang and Sellke \cite[arXiv version~2, Theorem~1, p.~4, and Corollary~1.8, p.~8]{hs2023}; see \cref{sel:published-field-threshold}. As the test direction is independent of the disorder, we do not need the adaptive disclosure argument that \cref{app:general-selection} uses for forgetting.

\paragraph{Setting.}
We work in the setting of \cref{app:general-selection}. We take dimension limits along the subsequence $(N_j)$ of \cref{thm:fisher-clock} and write $N\to\infty$ for $N=N_j\to\infty$. We keep every time horizon, branch number, and approximation parameter fixed while $N\to\infty$, and take long-time limits afterward.

Fix an integer $K\ge1$ and times $T>0$ and $L\ge0$. We run the diffusion \eqref{eq:hard-sde} up to time $T$, driven by the Brownian motion $(B_u)_{u\le T}$, which we call the trunk. From $X_T$, we run $K$ continuations of duration $L$, driven by Brownian motions independent of each other and of $(J,X_0,(B_u)_{u\le T})$. Let $Y_1,\dots,Y_K$ be their endpoints at time $T+L$, and let $m_{N,T}=K^{-1}\sum_{b=1}^KY_b$ be their mean, which also depends on $K$ and $L$. For each $b$, the Brownian path driving the $b$-th continuation, extended by the trunk on $[0,T]$, is a standard Brownian motion on $[0,T+L]$ independent of $(J,X_0)$. Since \eqref{eq:hard-sde} has a unique strong solution, the vector $(J,X_0,(B_u)_{u\le T},Y_b)$ has the same law as $(J,X_0,(B_u)_{u\le T},X_{T+L})$.

Let $Q_{K,L}$ and $h_{K,L}$ be given by \eqref{sel:fork-norm} and \eqref{sel:fork-center-energy}, with $e=e_\beta(\infty)$. When $q>0$, every sequence $T\to\infty$ realizes the stationary profile $(c,r)$ by \cref{sel:all-temperature}, and we may apply \cref{sel:finite-forks} along every such sequence. By \cref{lem:fork-parameter-limits}\ref{item:fork-K-limit}, we have, for every fixed $L\ge0$,
\[
 \lim_{K\to\infty}Q_{K,L}=c(2L)=d_L,\qquad \lim_{K\to\infty}h_{K,L}=h_L.
\]
We now show that the quantity $\delta_L$ defined at the start of \cref{resp:chapter} satisfies
\begin{equation}
 \delta_L=\Eth-\frac{h_L}{d_L^{p/2}}\ge0,
 \qquad \lim_{L\to\infty}\delta_L=0.
 \label{resp:p-center-deficit}
\end{equation}
The inequality is \eqref{eq:threshold-input-ceiling}; it also follows by letting $K\to\infty$ in \eqref{sel:fork-center-threshold}. The plateau is marginal by \cref{sel:positive-marginality}, so \eqref{sel:fork-saturation} applies, and it gives the limit, because $h_{K,L}/Q_{K,L}^{p/2}\to h_L/d_L^{p/2}$ as $K\to\infty$. For $p\in\{3,4\}$, the quantity $\delta_L$ is the same as in \eqref{sel:center-deficit}.

\paragraph{Two lemmas.}
The first lemma converts the absolute-overlap cost of \cref{sel:absolute-field-cost} into a bound on the overlap with a Gaussian vector from an independent seed. We write $B_N=\{x\in\R^N:\norm x\le\sqrt N\}$. Recall that $\bd=\bd(p)$ is a constant of the model, defined in \cref{sec:intro-picture}; it does not depend on the inverse temperature $\beta$.

\begin{lemma}[Gaussian seed overlap has a threshold cost]
\label{resp:seed-cost}
Fix $p\ge3$ and $0\le L_0<\infty$. For each $N$, let $\zeta_N$ be a random element independent of the coefficient tensor $J$, let $G_N$ be a standard Gaussian vector in $\R^N$ that is a measurable function of $\zeta_N$, and let $A_N=A_N(J,\zeta_N)\in B_N$ be measurable. Suppose that, for almost every value of $\zeta_N$, the map $J\mapsto A_N(J,\zeta_N)$ is globally $L_0$-Lipschitz, for the Euclidean distances on the ordered coefficient tensors and on $\R^N$. Let $\delta\ge0$, and suppose that
\[
 \liminf_{N\to\infty}\frac{\E H_N(A_N)}N\ge\Eth-\delta.
\]
Then
\begin{equation}
 \limsup_{N\to\infty}\E\Bigl|\frac{\ip{A_N}{G_N}}N\Bigr|\le\bd^{-1/p}\,\delta^{1/p}.
 \label{resp:seed-cost-bound}
\end{equation}
The constant $\bd^{-1/p}$ depends neither on $L_0$ nor on the law of $\zeta_N$. The same holds when $N$ runs through a subsequence.
\end{lemma}

The algorithm $A_N$ may depend on $G_N$; only $J$ must be independent of the seed. If $A_N$ also depends on parameters held fixed while $N\to\infty$, the lemma applies at each of their values, and its conclusion passes to limits in these parameters taken afterward.

\begin{proof}
Let $w_N=\sqrt N\,G_N/\norm{G_N}$ if $G_N\ne0$, and $w_N=\sqrt N\,e_1$ otherwise. Then $w_N$ is a measurable function of $\zeta_N$ with $\norm{w_N}=\sqrt N$. Since $\norm{A_N}\le\sqrt N$, we have
\[
 \frac{|\ip{A_N}{G_N}-\ip{A_N}{w_N}|}N\le\frac{\norm{G_N-w_N}}{\sqrt N}=\Bigl|\frac{\norm{G_N}}{\sqrt N}-1\Bigr|.
\]
For $x\ge0$, we have $|x-1|\le|x^2-1|$. The variable $\norm{G_N}^2/N$ has mean one and variance $2/N$, and the expectation of the right side is at most $\E|\norm{G_N}^2/N-1|\le\sqrt{2/N}$, which tends to zero.

We use \cref{sel:absolute-field-cost} conditionally on $\zeta_N$, with direction $w_N$, which is $\zeta_N$-measurable, and Lipschitz constant $L_0$. Writing $a_N=\ip{A_N}{w_N}/N$, we have
\[
 \limsup_{N\to\infty}\Bigl\{\frac{\E H_N(A_N)}N+\bd\,\E|a_N|^p\Bigr\}\le\Eth.
\]
Since $\limsup_N(x_N+y_N)\ge\liminf_Nx_N+\limsup_Ny_N$, the hypothesis on the energy implies that $\bd\limsup_N\E|a_N|^p\le\delta$. Hölder's inequality gives $\E|a_N|\le(\E|a_N|^p)^{1/p}$, and hence $\limsup_N\E|a_N|\le(\delta/\bd)^{1/p}$. Combining this with the first paragraph proves \eqref{resp:seed-cost-bound}. The constants in \cref{sel:absolute-field-cost} depend neither on the seed nor on the direction.
\end{proof}

The second lemma identifies the overlap of the mean $m_{N,T}$ with a Gaussian vector built from the trunk. It is a weighted form of the Brownian response identity of \cref{lem:brown-force-response} and \cref{prop:hard-bridge}.

\begin{lemma}[A weighted Brownian response identity]
\label{resp:weighted-brownian}
Fix $K\ge1$, $T>0$, and $L\ge0$, and let $k:[0,T]\to\R$ be a continuous deterministic function with $\int_0^Tk(u)^2\,\dd u=1$. Let $G_N=\int_0^Tk(u)\,\dd B_u$, where $B$ is the trunk Brownian motion. Then $G_N$ is a standard Gaussian vector in $\R^N$, independent of the disorder, and
\begin{equation}
 \lim_{N\to\infty}\E\Bigl[\frac{\ip{m_{N,T}}{G_N}}N\Bigr]
 =\sqrt2\int_0^Tk(u)R(T+L,u)\,\dd u.
 \label{resp:weighted-brownian-identity}
\end{equation}
\end{lemma}

\begin{proof}
The coordinates of $G_N$ are Wiener integrals of the independent coordinates of $B$. They are independent centered Gaussian variables with variance $\int_0^Tk^2=1$, and they are independent of $(J,X_0)$ because $B$ is.

Let $t=T+L$. Since $G_N$ is a function of the trunk, the equality of laws noted above gives $\E\ip{Y_b}{G_N}=\E\ip{X_t}{G_N}$ for each $b$. Averaging over $b$, we have $\E\ip{m_{N,T}}{G_N}=\E\ip{X_t}{G_N}$. Suppose first that $k=\sum_{i=1}^nk_i\mathbf1_{(u_{i-1},u_i]}$ is a step function, with $0=u_0<\dots<u_n=T$. Then
\[
 \frac{\E\ip{X_t}{G_N}}N=\sum_{i=1}^nk_i\,\frac{\E\ip{X_t}{B_{u_i}-B_{u_{i-1}}}}N
 =\sqrt2\sum_{i=1}^nk_i\bigl[\E\Gamma_N(t,u_i)-\E\Gamma_N(t,u_{i-1})\bigr],
\]
with $\Gamma_N$ as in \cref{prop:hard-bridge}. That proposition states that $\E\Gamma_N(t,u)\to\int_0^uR(t,w)\,\dd w$, and the right side converges to $\sqrt2\int_0^Tk(u)R(t,u)\,\dd u$.

For a general continuous $k$ and $\varepsilon>0$, we choose a step function $k_\varepsilon$ with $\sup_{[0,T]}|k-k_\varepsilon|\le\varepsilon$. We have $\norm{X_t}^2=N$ and, by the Itô isometry, $\E\norm{\int_0^T(k-k_\varepsilon)\,\dd B}^2=N\int_0^T(k-k_\varepsilon)^2$. Using the Cauchy--Schwarz inequality, we find
\[
 \frac1N\Bigl|\E\Bigl\langle X_t,\int_0^T(k-k_\varepsilon)\,\dd B\Bigr\rangle\Bigr|
 \le\Bigl(\int_0^T(k-k_\varepsilon)^2\Bigr)^{1/2}\le\sqrt T\,\varepsilon,
\]
uniformly in $N$. Also $|\sqrt2\int_0^T(k-k_\varepsilon)R(t,\cdot)|\le\sqrt2\,T\varepsilon$, since $0\le R\le1$. Using the step-function case for $k_\varepsilon$, we obtain
\[
 \limsup_{N\to\infty}\Bigl|\frac{\E\ip{X_t}{G_N}}N-\sqrt2\int_0^Tk(u)R(t,u)\,\dd u\Bigr|\le(\sqrt T+\sqrt2\,T)\varepsilon.
\]
Letting $\varepsilon\downarrow0$ proves \eqref{resp:weighted-brownian-identity}.
\end{proof}

The Brownian observable goes back to Ben Arous, Dembo, and Guionnet \cite[arXiv version~1, equation~(1.9), p.~4, and Theorem~1.2, p.~5]{bag2006}, and its hard-sphere limit is \cite[arXiv version~5, Definition~3, p.~22, and Theorem~4.3, p.~23]{sellke2024}, which underlies \cref{prop:hard-bridge}. By \cref{lem:brown-force-response}, the left side of \eqref{resp:weighted-brownian-identity} at finite $N$ equals $\sqrt2\int_0^Tk(u)\,\bar r_N(T+L,u)\,\dd u$ when $k$ is a step function; this is the Cameron--Martin identity for a deterministic force with profile $k$. We need only its limit, which \cref{prop:hard-bridge} provides without a pointwise limit of the finite-dimensional response.

\paragraph{The remote response.}
We now prove the remaining case of \eqref{eq:resp-remote-vanishing}.

\begin{proposition}[Remote squared response for general pure models]\label{prop:remote-response-general}
Fix $p\ge3$ and $\beta>0$, and suppose that the plateau $q$ is positive. Then, for every $L>0$,
\begin{equation}\label{eq:resp-general-remote-bound}
 \limsup_{t\to\infty}\int_0^{t-L}R(t,u)^2\,\dd u\le\frac{d_L}2\,\bd^{-2/p}\,\delta_L^{2/p}.
\end{equation}
In particular, \eqref{eq:resp-remote-vanishing} holds.
\end{proposition}

The idea is as follows. We choose the weight $k$ in \cref{resp:weighted-brownian} proportional to the remote response row $u\mapsto R(T+L,u)$ on $[0,T]$. Then the overlap of $m_{N,T}$ with $G_N$ is $\sqrt{2S_2(T+L,L)}$ in the limit. After normalization by $\sqrt{Q_{K,L}}$, the mean $m_{N,T}$ has energy close to $\Eth$, and \cref{resp:seed-cost} forces this overlap to be small. The coefficient-Lipschitz approximations of \cref{sel:finite-forks} make \cref{resp:seed-cost} applicable.

\begin{proof}
Fix $L>0$ and $K\ge1$, and abbreviate $Q=Q_{K,L}$ and $\delta_{K,L}=\Eth-h_{K,L}/Q^{p/2}$. Recall from \eqref{sel:fork-center-threshold} that $\delta_{K,L}\ge0$, and note that $Q\ge c(2L)\ge q>0$. The proof has four steps. We take the limits in the following order:
\begin{enumerate}[label=(\roman*)]
\item $N=N_j\to\infty$, with $T$, $K$, $L$, and the approximation parameter $\varepsilon$ fixed;
\item $T\to\infty$, along an arbitrary sequence;
\item $\varepsilon\downarrow0$;
\item $K\to\infty$.
\end{enumerate}

\paragraph{Step 1: the test direction.}
For $T>0$, the response is positive by \eqref{eq:resp-response-lower}, and $S_2(T+L,L)=\int_0^TR(T+L,u)^2\,\dd u>0$. We define
\[
 k_{T,L}(u)=\frac{R(T+L,u)}{S_2(T+L,L)^{1/2}},\qquad 0\le u\le T,
 \qquad G_N=\int_0^Tk_{T,L}(u)\,\dd B_u.
\]
The weight $k_{T,L}$ is continuous and deterministic: it depends on the CK solution and on the fixed times $T$ and $L$, not on the disorder or the dimension. It satisfies $\int_0^Tk_{T,L}^2=1$. \Cref{resp:weighted-brownian} implies that
\begin{equation}\label{eq:resp-test-overlap}
 \lim_{N\to\infty}\E\Bigl[\frac{\ip{m_{N,T}}{G_N}}N\Bigr]
 =\sqrt2\,\frac{\int_0^TR(T+L,u)^2\,\dd u}{S_2(T+L,L)^{1/2}}
 =\sqrt{2S_2(T+L,L)}.
\end{equation}

\paragraph{Step 2: the normalized mean and its Lipschitz approximation.}
Let $\Pi_{B_N}$ be the metric projection onto $B_N$, and let $\bar m_{N,T}=\Pi_{B_N}(m_{N,T}/\sqrt Q)$. For $\varepsilon>0$, let $\mathcal A^\varepsilon_{N,T}$ be the maps of \cref{sel:finite-forks}, so that
\[
 \limsup_{N\to\infty}\frac1N\E\norm{\mathcal A^\varepsilon_{N,T}-\bar m_{N,T}}^2\le\varepsilon^2,
\]
as in \eqref{sel:fork-lipschitz}. These maps take values in $B_N$. Conditionally on the initial point and on all Brownian paths, they are globally Lipschitz in $J$, with a constant that depends neither on $N$ nor on the seed. The vector $G_N$ is a function of the trunk Brownian path. It is therefore measurable with respect to this seed, which is independent of $J$.

We next compare energies. The vectors $m_{N,T}/\sqrt Q$, $\bar m_{N,T}$, and $\mathcal A^\varepsilon_{N,T}$ have norm at most $\sqrt{N/q}$. By homogeneity and \eqref{sel:tensor-ui}, we have $\norm{\nabla H_N(x)}\le pQ_Nq^{-(p-1)/2}\sqrt N$ when $\norm x\le\sqrt{N/q}$, where $Q_N$ is the tensor norm \eqref{eq:tensor-norm}. Integrating along segments and applying the Cauchy--Schwarz inequality, we find that, for random vectors $x,y$ of norm at most $\sqrt{N/q}$,
\[
 \Bigl|\frac{\E H_N(x)-\E H_N(y)}N\Bigr|\le D_q\Bigl(\frac{\E\norm{x-y}^2}N\Bigr)^{1/2},
 \qquad D_q=pq^{-(p-1)/2}\sup_N(\E Q_N^2)^{1/2},
\]
and $D_q<\infty$ by \cref{lem:landscape-moments}. We set
\[
 \gamma_T=\limsup_{N\to\infty}\Bigl|\frac{\E H_N(m_{N,T})}N-h_{K,L}\Bigr|,
 \qquad
 \kappa_T=\limsup_{N\to\infty}\Bigl(\frac{\E\norm{\bar m_{N,T}-m_{N,T}/\sqrt Q}^2}N\Bigr)^{1/2}.
\]
Note that $H_N(m_{N,T}/\sqrt Q)=Q^{-p/2}H_N(m_{N,T})$ by homogeneity. Comparing energies through the intermediate vector $\bar m_{N,T}$, we obtain
\begin{equation}\label{eq:resp-algorithm-energy}
\begin{aligned}
 \liminf_{N\to\infty}\frac{\E H_N(\mathcal A^\varepsilon_{N,T})}N
 &\ge Q^{-p/2}(h_{K,L}-\gamma_T)-D_q(\varepsilon+\kappa_T)\\
 &=\Eth-\delta_{K,L}-Q^{-p/2}\gamma_T-D_q(\varepsilon+\kappa_T).
\end{aligned}
\end{equation}

\paragraph{Step 3: the seed cost.}
By Step 2, the maps $\mathcal A^\varepsilon_{N,T}$ satisfy the hypotheses of \cref{resp:seed-cost}. The seed is formed by the initial point and all Brownian paths, the Gaussian vector is the vector $G_N$ of Step 1, and $\delta=\delta_{K,L}+Q^{-p/2}\gamma_T+D_q(\varepsilon+\kappa_T)$ by \eqref{eq:resp-algorithm-energy}. \Cref{resp:seed-cost} implies that
\[
 \limsup_{N\to\infty}\E\Bigl|\frac{\ip{\mathcal A^\varepsilon_{N,T}}{G_N}}N\Bigr|
 \le\bd^{-1/p}\bigl(\delta_{K,L}+Q^{-p/2}\gamma_T+D_q(\varepsilon+\kappa_T)\bigr)^{1/p}.
\]
To return to $m_{N,T}$, we use $\E\norm{G_N}^2=N$ and the Cauchy--Schwarz inequality; the fact that $G_N$ and the approximation share the trunk seed plays no role here. We have
\[
 \Bigl|\E\frac{\ip{m_{N,T}/\sqrt Q-\mathcal A^\varepsilon_{N,T}}{G_N}}N\Bigr|
 \le\Bigl(\frac{\E\norm{m_{N,T}/\sqrt Q-\bar m_{N,T}}^2}N\Bigr)^{1/2}
 +\Bigl(\frac{\E\norm{\bar m_{N,T}-\mathcal A^\varepsilon_{N,T}}^2}N\Bigr)^{1/2}.
\]
Combining the last two displays with \eqref{eq:resp-test-overlap}, we find that, for every $T>0$,
\begin{equation}\label{eq:resp-fixed-T-bound}
 \sqrt{\frac{2S_2(T+L,L)}Q}
 \le\bd^{-1/p}\bigl(\delta_{K,L}+Q^{-p/2}\gamma_T+D_q(\varepsilon+\kappa_T)\bigr)^{1/p}+\varepsilon+\kappa_T.
\end{equation}

\paragraph{Step 4: the long-time limits.}
Let $T_n\to\infty$ be an arbitrary sequence. Using \eqref{sel:fork-center-energy-limit}, we have $\gamma_{T_n}\to0$. For $\kappa_T$, note that $\norm{\Pi_{B_N}(x)-x}=(\norm x-\sqrt N)_+$, and that $|\sqrt{y/Q}-1|\le|y-Q|/Q$ for $y\ge0$. With $y=\norm{m_{N,T}}^2/N$, we find $\kappa_T^2\le\limsup_N\E|\norm{m_{N,T}}^2/N-Q|^2/Q^2$, and \eqref{sel:fork-center-concentration} implies that $\kappa_{T_n}\to0$. The left side of \eqref{eq:resp-fixed-T-bound} does not depend on $\varepsilon$. Taking $T=T_n\to\infty$ in \eqref{eq:resp-fixed-T-bound} and then $\varepsilon\downarrow0$, we obtain
\[
 \limsup_{n\to\infty}S_2(T_n+L,L)\le\frac Q2\,\bd^{-2/p}\,\delta_{K,L}^{2/p}.
\]
Since the sequence $T_n$ was arbitrary, the same bound holds for $\limsup_{t\to\infty}S_2(t,L)$. This upper limit does not depend on $K$. As $K\to\infty$, we have $Q_{K,L}\to d_L$ and $\delta_{K,L}\to\delta_L$, by \cref{lem:fork-parameter-limits}\ref{item:fork-K-limit}. This proves \eqref{eq:resp-general-remote-bound}. As $d_L\le1$ and $\delta_L\to0$ by \eqref{resp:p-center-deficit}, the right side of \eqref{eq:resp-general-remote-bound} tends to zero as $L\to\infty$, which is \eqref{eq:resp-remote-vanishing}.
\end{proof}

Together with Steps~1, 2, and~4 of the proof of \cref{resp:strong-row}, \cref{prop:remote-response-general} proves strong $L^2$ convergence of the response rows, the exact limits of the squared tails, and $\chi(t,0)=o(\sqrt t)$ for every pure model with $p\ge3$.
